\section{Expansion as Projective History Branching}
\label{sec:history-branching-expansion}

The preceding sections describe the late-time phenomenological sector of the model.
We now isolate a distinct structural question: whether the accelerated component of cosmic expansion can be
associated with the growth of admissible projective histories rather than with a material dark-energy fluid.

The relevant result is not an exponential growth of Heisenberg balls.
Indeed, the spectral admissibility cascade is carried by the Heisenberg group, whose endpoint balls have
polynomial Bass--Guivarc'h growth of degree four before finite-$q$ saturation.
Endpoint counting therefore cannot support an asymptotic de Sitter-like component.
The surviving exponential object is instead the number of projectively distinguishable
histories~\cite{Beau2026tb}.
For the canonical O12-compatible residual channel, the trajectory-branching note proves that endpoint-based
distinguishability collapses, whereas D4 projective distinguishability yields the exact Pell count
\[
N_{\mathrm{dist}}(n)=N_b(n),
\qquad
N_b(n)=2N_b(n-1)+N_b(n-2),
\]
with asymptotic rate
\[
h_b=\lim_{n\to\infty}\frac{1}{n}\log N_b(n)
=\log(1+\sqrt{2}).
\]
This relocates the exponential component of the two-regime expansion model of the spectral relaxation
paper~\cite{Beau2026a5} from state volume to history branching.

\begin{hypothesis}[H-dict: rank--time dictionary]
\label{hyp:H-dict}
The coarse-grained homogeneous cosmological update identifies one admissible cascade rank with one projective
clock tick of duration $H^{-1}$.
Equivalently,
\[
dt = H^{-1}dn,
\qquad
dn = H\,dt = d\log a.
\]
Thus, up to an additive constant, the cascade rank is the number of e-folds,
\[
n=\log a.
\]
\end{hypothesis}

Hypothesis~\ref{hyp:H-dict} is a dictionary between the temporal-ordering sector and the homogeneous
cosmological sector.
It is not an additional expansion law.
It fixes how the already-defined projective rank is read by the effective FRW observer.
Its two ingredients are already separately established: the temporal-ordering sector identifies one
admissible cascade rank with one intrinsic tick~\cite{Beau2026tp}, and the operational closure of
Appendix~\ref{app:cosmo_resolution} identifies $H^{-1}$ as the unique intrinsic update timescale of the
homogeneous regime.
What Hypothesis~\ref{hyp:H-dict} adds --- and what keeps it a hypothesis rather than a theorem --- is the
identification of the discrete projective rank of the substrate cascade with the coarse-grained homogeneous
update step, a bridge between sub-programmes.

Under Hypothesis~\ref{hyp:H-dict}, the canonical projective history count scales as
\[
N_{\mathrm{dist}}(a)\sim a^{h_b},
\qquad
h_b=\log(1+\sqrt{2})\simeq 0.881.
\]
This immediately rules out the naive volumetric interpretation $N_{\mathrm{dist}}\propto a^3$: the exponent
mismatch ($0.881$ against $3$) makes the identification quantitatively impossible, independently of the
structural endpoint no-go.
The branching count is not the comoving volume of space.
It is the entropy-bearing count of distinguishable reprojection histories.

The de Sitter-like reading is therefore structural.
The accelerated component is not sourced by a material fluid with negative pressure.
It reflects the persistence of a positive projective history-branching rate: the admissible space of future
reprojections does not terminate under the Heisenberg transfer.

Two clarifications protect this statement.
First, under Hypothesis~\ref{hyp:H-dict} the exponential growth in rank is a power law in the scale factor by
construction; the branching result therefore does not by itself determine $H(t)$, and in particular does not
derive the acceleration of the scale factor.
What it establishes is that the projective sector sustains, rather than obstructs, unbounded e-folding, and
that it assigns to the accelerated phase a definite entropy carrier, with history-entropy production rate
$\dot S_{\mathrm{hist}} = h_b\,H$.
Second, the value of the asymptotic Hubble scale $H_\infty$ is not fixed by the Pell rate alone.
It requires the normalization of projective rank against the physical homogeneous clock and is left as an
open normalization problem.
The named corpus route towards such a normalization is the variational principle of the spectral-equilibrium
sector, where the Einstein response arises as an Euler--Lagrange condition of a renormalized spectral
entropy~\cite{Beau2026Thermodynamics}; connecting the history-branching entropy to that variational object is
an open front, not used here.

The canonical channel is rigid and gives a constant branching rate.
A time-dependent effective equation of state, if present, must come from the compression of full-history
profiles rather than from the $b$-shadow channel.
The full-history Gabor channel provides the current candidate for such a compression
profile~\cite{Beau2026tb}, but its canonical cosmological interpretation remains open.

\subsection{Large-angle imprint of the pre-branching regime}
\label{subsec:low-ell-prebranching}

The projective history-branching result also clarifies a possible origin of the low-$\ell$ anomaly discussed
in the earlier white-paper formulation~\cite{Cosmochrony}.
In the canonical D4 channel, the first ranks are projectively degenerate: the number of distinguishable
histories is small before the Pell branching regime becomes effective.
Under Hypothesis~\ref{hyp:H-dict}, the smallest cascade ranks correspond to the largest observable scales.
A suppression of distinguishable projective histories at $n\leq 2$ is therefore a natural candidate mechanism
for a suppression of large-angle variance.

This statement should not be confused with a derivation of the observed low-$\ell$ spectrum.
The earlier ratio $\lambda_2/\lambda_1\simeq 8/3$ belongs to the $S^3$/Hopf spectral relaxation experiments
of the white paper and must be re-audited in the Heisenberg D4 setting.
The present result supplies the missing structural mechanism: a pre-branching window in which projectively
distinguishable histories have not yet reached their asymptotic entropy rate.
The quantitative link to CMB multipoles remains open.
